\section{Discussion}\label{sec:discussion}

We found PLS an increasingly popular area of study, but researchers primarily rely on a handful of traditional metrics for evaluation. However, we found that traditional metrics are imperfect measures of readability and LM evaluators can draw significantly different, and more accurate, conclusions about PLS datasets than when using FKGL, the most common metric.



\subsection{Why traditional readability metrics are insufficient measures of readability}\label{subsec:discussion-inconsistency}
% \daniel{if this does not need to be a "subsection" consider turning it into a paragraph and it will save you some space. }
We consider 2 explanations for the poor correlation of readability metrics with human judgments: definitional inconsistency or measurement error. Definitional inconsistency means that the definition of ``readable,'' as measured by the metrics, differs from the definition of ``readable,'' as considered by human judges. Measurement error means that, even if we have the correct definition, we are not measuring readability properly.
%This distinction is important because it allows us to understand the source of the low performance and design appropriate solutions.
We argue that there is evidence for both problems.

On definitional inconsistency, the majority of readability metrics originated in the education domain. Traditional readability metrics typically define a ``readable'' text as one with an appropriate text complexity for the number of years of education (i.e., a text has a US 9th grade reading level)~\parencite{Gunning1968TheTO,Coleman1975ACR,flesch1952simplification}. In contrast, the field of PLS typically defines a ``readable'' text as one that gives a non-expert, adult reader an overall understanding of the source article. These different definitions have different implications for the resulting text. If optimizing for education-appropriate text complexity, we can measure the complexity of the vocabulary or sentences. However, using the PLS definition of readability, we should measure features such as whether the text includes explanations of technical terms or how much background is required to understand the concepts.

Traditional readability metrics also suffer from measurement error. Even if we assume a consistent definition, traditional metrics do not properly measure readability. They measure lexical properties, such as number of syllables in a word, which penalizes summaries for using clearly defined technical terms. Traditional metrics also do not measure deeper features that contribute to readability, such as how much background is required to understand the text. In \Cref{subsec:results-rqthree}, we show that LMs are better able to reason over these more complex attributes. 

\Cref{tab:disagreement-example} shows examples in which the LM evaluator and FKGL disagree on the readability. FKGL rates a summary from the SJK dataset as having a graduate-college reading level, while the LM rates it as highly readable (\Cref{subtab:ex-disaggreement-sjk}). Although the summary explains the concepts well, long words such as ``harvesting'' and ``electricity'' likely cause FKGL to rate the summary as less readable. \Cref{subtab:ex-disaggreement-plos} has a Pubmed example, which the LM rates as having low readability, while FKGL assigns the summary a 10th grade reading level. This example contains many short words, such as ``GCDC'' and ``SPE'', which are favored by FKGL. Although short, these technical words are not well defined. For example, the  ``GCDC'' is defined as ``glycochenodeoxycholic acid,'' but  is not otherwise explained. In general, we  notice that FKGL favors acronyms, which are often present in technical text.

\subsection{Recommendations and Future Directions}\label{subsec:recommendations}
We find that many traditional readability metrics have poor correlation with human judgments and that LMs provide better judgments. However, LM-evaluators are an imperfect solution since they are subject to bias and a lack of interpretability~\cite{Liu2023LLMsAN, Wang2023LargeLM, Shen2023LargeLM, Stureborg2024LargeLM}. Therefore, we recommend a multi-faceted evaluation of PLS that uses a combination of traditional readability metrics and LM evaluators. Specifically, we recommend using DCRS and CLI, which have the highest correlation with human judgments. We recommend discontinuing use of FKGL for PLS, the current most popular metric, due to low correlation with human judgment. We recommend using LMs as additional metrics, especially for more qualitative evaluations, such as the keyword analysis conducted in \Cref{subsec:results-rqthree}. These types of analyses give a more holistic view of the benefits and downsides of datasets and methods. Finally, we recommend that PLS research use datasets with higher readability scores (\Cref{subsec:results-rqthree}), such as CDSR and SciNews. We recommend that PLOS and CELLS be considered general scientific summarization datasets and not plain language datasets. This recommendation is particularly impactful as PLOS has been used in every year the shared task BioLaySumm has occurred~\cite{goldsack-etal-2024-biolaysumm,biolaysumm-2023-overview}.

Future work should focus on constructing metrics that better align with human judgments of readability in both definition and measurement (\Cref{subsec:discussion-inconsistency}). We show that LMs are promising and worthy of future work that can decrease bias and improve interpretability. Dataset collection should focus on collecting highly readable summaries and consider deeper attributes of readable summaries, such as explanations of technical concepts. 